\documentclass[aspectratio=169, 10pt]{beamer}
\usepackage{sunass_estilo}

\title{Módulo III: Competencia Perfecta}
\subtitle{Guía Teórica de Repaso Rápido · PEU SUNASS}
\author[PEU SUNASS · Economía]{Preparación Examen de Admisión PEU SUNASS}
\institute{Basado en el Temario Oficial de Economía · Microeconomía}
\date[2026]{2026}

\begin{document}

% ==============================================================================
% SLIDE 1: PORTADA
% ==============================================================================
\begin{frame}[plain]
  \titlepage
\end{frame}

% ==============================================================================
% SLIDE 2: SUPUESTOS Y MAPA CONCEPTUAL
% ==============================================================================
\begin{frame}{Mapa Conceptual y Supuestos de la Competencia Perfecta}
  \vspace{0.15cm}
  \begin{columns}[T]
    \begin{column}{0.32\textwidth}
      \tarjeta{azulPrincipal}{1. Supuestos Estructurales}{
        • \textbf{Atomicidad:} muchos compradores y vendedores pequeños (precio-aceptantes).\\
        • \textbf{Homogeneidad:} producto idéntico (sustituto perfecto).\\
        • \textbf{Información perfecta} y cero costos de transacción.
      }
    \end{column}
    \begin{column}{0.32\textwidth}
      \tarjeta{naranjaAcento}{2. Equilibrio CP y LP}{
        • \textbf{Corto plazo:} número fijo de firmas \(N\); puede haber \(\pi > 0\) o \(\pi < 0\).\\
        • \textbf{Largo plazo:} libre entrada y salida \(\implies \pi = 0\) en equilibrio.
      }
    \end{column}
    \begin{column}{0.32\textwidth}
      \tarjeta{verdeAcento}{3. Bienestar e Intervención}{
        • Maximiza \(W = EC + EP\) (eficiencia asignativa y productiva).\\
        • Análisis de impuestos, subsidios, precios tope y aranceles (\(PIE\)).
      }
    \end{column}
  \end{columns}

  \vspace{0.2cm}
  \recuadroclave{
    \small\textbf{Consecuencia clave:} Cada firma enfrenta una curva de demanda individual \textbf{perfectamente elástica} (\(|\varepsilon_d| = \infty\)) al precio de mercado: \(P = IMe = IMg\).
  }
  \mensajefinal{En competencia perfecta ninguna empresa tiene poder de mercado: el Índice de Lerner es \(L = \frac{P-CMg}{P} = 0\).}
\end{frame}

% ==============================================================================
% SLIDE 3: MAXIMIZACIÓN DEL BENEFICIO Y OFERTA DE CORTO PLAZO
% ==============================================================================
\begin{frame}{1. Maximización del Beneficio y Oferta de la Industria (CP)}
  \begin{columns}[T]
    \begin{column}{0.53\textwidth}
      {\small\textbf{Agregación Horizontal de Ofertas Individuales}}
      \vspace{0.05cm}
      \begin{tikzpicture}[scale=0.76]
        % Panel Izq: Firma
        \begin{axis}[
          axis lines=left,
          xlabel={$q_i$}, ylabel={$P$},
          xmin=0, xmax=5.5, ymin=0, ymax=9,
          xtick={3.5}, xticklabels={$q_i^*$},
          ytick={2.5, 6}, yticklabels={$\min CVMe$, $P^*$},
          width=4.2cm, height=5.0cm,
          name=firma
        ]
          \addplot[azulPrincipal, very thick, domain=1.5:4.1, samples=30] {2.5 + 0.875*(x-1.5)^2};
          \addplot[grisTexto, dashed, domain=0.3:1.5, samples=20] {2.5 + 0.875*(x-1.5)^2};
          \draw[rojoAcento, very thick] (axis cs:0, 6) -- (axis cs:5.0, 6);
          \draw[dotted, grisTexto] (axis cs:3.5, 0) -- (axis cs:3.5, 6);
          \addplot[only marks, mark=*, mark size=2pt, rojoAcento] coordinates {(3.5, 6)};
          \node[azulPrincipal, font=\tiny, anchor=east] at (axis cs:3.9, 8.2) {$CMg_i$};
        \end{axis}

        % Panel Der: Industria (N firmas)
        \begin{axis}[
          at={(firma.east)}, anchor=west, xshift=1.0cm,
          axis lines=left,
          xlabel={$Q = \sum q_i$}, ylabel={$P$},
          xmin=0, xmax=8, ymin=0, ymax=9,
          xtick={4.5}, xticklabels={$Q^*$},
          ytick={2.5, 6}, yticklabels={$\min CVMe$, $P^*$},
          width=4.4cm, height=5.0cm
        ]
          \addplot[verdeAcento, very thick, domain=1.5:5.5, samples=30] {2.5 + 0.3889*(x-1.5)^2};
          \addplot[naranjaAcento, very thick, domain=0.8:7.4] {9.2 - 0.7111*x};
          \draw[dotted, grisTexto] (axis cs:0, 6) -- (axis cs:4.5, 6) -- (axis cs:4.5, 0);
          \addplot[only marks, mark=*, mark size=2pt, verdeAcento] coordinates {(4.5, 6)};
          \node[verdeAcento, font=\tiny, anchor=east] at (axis cs:5.3, 8.2) {$S_{CP}$};
          \node[naranjaAcento, font=\tiny] at (axis cs:6.8, 3.8) {$D$};
        \end{axis}
      \end{tikzpicture}
    \end{column}

    \begin{column}{0.44\textwidth}
      {\scriptsize\color{grisTexto}Condición de Primer y Segundo Orden}\\[1pt]
      {\large\color{azulPrincipal}\textbf{\(P = CMg(q_i) \quad \text{y} \quad \frac{d\,CMg}{dq_i} > 0\)}}\\[5pt]
      {\scriptsize
        • La firma produce en el tramo \textbf{creciente} del \(CMg\) por encima de \(\min CVMe\).\\
        • Si \(P = CMg\) se cumple en el tramo decreciente, es un mínimo de beneficio.\\[5pt]
      }
      \indicador{Oferta de la Industria (CP)}{verdeAcento}{\(Q^s(P) = \sum_{i=1}^N q_i^s(P)\)}
      {\scriptsize
        • Es la \textbf{suma horizontal} de las ofertas individuales.\\
        • Si los precios de factores suben al expandirse toda la industria, \(S_{CP}\) es más inelástica que la suma simple.
      }
    \end{column}
  \end{columns}

  \mensajefinal{Un impuesto de suma alzada (fijo) o un alza del costo fijo NO desplaza \(S_{CP}\); solo reduce \(\pi\) en el corto plazo.}
\end{frame}

% ==============================================================================
% SLIDE 4: EQUILIBRIO A LARGO PLAZO (LIBRE ENTRADA Y SALIDA)
% ==============================================================================
\begin{frame}{2. Equilibrio Competitivo a Largo Plazo (\(\pi = 0\))}
  \begin{columns}[T]
    \begin{column}{0.53\textwidth}
      {\small\textbf{Ajuste de Entrada/Salida hasta Beneficio Nulo}}\\[3pt]
      \begin{tikzpicture}[scale=0.78]
        \begin{axis}[
          axis lines=left,
          xlabel={Producción ($q_i$)}, ylabel={Precio / Costo},
          xmin=0, xmax=8, ymin=0, ymax=10,
          xtick={4.0}, xticklabels={$q_{EME}$},
          ytick={3.5}, yticklabels={$P_{LP}^*$},
          width=8.2cm, height=5.2cm
        ]
          % CMe_LP: 3.5 + 0.35*(x-4)^2
          \addplot[azulPrincipal, very thick, domain=1.0:7.0, samples=50] {3.5 + 0.35*(x-4)^2};
          \node[azulPrincipal, font=\scriptsize, anchor=west] at (axis cs:6.4, 6.0) {$CMe_{LP}$};
          % CMg_LP exacto: d[x*CMe]/dx = 3.5 + 0.35*(x-4)^2 + 0.7*x*(x-4) -> corta desde abajo en (4, 3.5)
          \addplot[rojoAcento, very thick, domain=2.1:5.3, samples=40] {3.5 + 0.35*(x-4)^2 + 0.7*x*(x-4)};
          \node[rojoAcento, font=\scriptsize, anchor=east] at (axis cs:5.15, 8.6) {$CMg_{LP}$};
          % Precio LP = min CMe_LP
          \draw[verdeAcento, very thick] (axis cs:0, 3.5) -- (axis cs:7.5, 3.5);
          \node[verdeAcento, font=\scriptsize, anchor=south west] at (axis cs:0.3, 3.6) {\(P = IMg\)};
          % Punto optimo
          \addplot[only marks, mark=*, mark size=2.6pt, verdeAcento] coordinates {(4.0, 3.5)};
          \draw[dotted, grisTexto] (axis cs:4.0, 0) -- (axis cs:4.0, 3.5);
        \end{axis}
      \end{tikzpicture}
    \end{column}

    \begin{column}{0.44\textwidth}
      \indicador{Condición Triple de Equilibrio LP}{verdeAcento}{\(P^* = CMg_{LP} = \min CMe_{LP}\)}
      {\scriptsize
        • \textbf{Eficiencia asignativa:} \(P = CMg\) (no hay peso muerto).\\
        • \textbf{Eficiencia productiva:} se produce en la Escala Mínima Eficiente (\(\min CMe_{LP}\)).\\
        • \textbf{Beneficio económico nulo (\(\pi = 0\)):} no hay incentivos para entrar ni salir. Incluye el \textbf{beneficio normal} (costo de oportunidad del capital).\\[5pt]
      }
      \indicador{Número de Firmas en Equilibrio}{azulPrincipal}{\(N^* = \frac{Q^d(P_{LP}^*)}{q_{EME}}\)}
      {\scriptsize
        Si aumenta la demanda de mercado, a CP sube \(P\) y \(\pi>0\), entran nuevas firmas, \(S_{CP}\) se desplaza a la derecha y \(P\) retorna a \(\min CMe_{LP}\).
      }
    \end{column}
  \end{columns}

  \mensajefinal{En el equilibrio a largo plazo, el beneficio económico es cero (\(\pi=0\)), pero el beneficio contable es positivo.}
\end{frame}

% ==============================================================================
% SLIDE 5: OFERTA DE LA INDUSTRIA A LARGO PLAZO (3 CASOS CLÁSICOS)
% ==============================================================================
\begin{frame}{3. Oferta de la Industria a Largo Plazo (\(S_{LP}\))}
  \vspace{0.1cm}
  \begin{columns}[T]
    \begin{column}{0.32\textwidth}
      \tarjeta{azulPrincipal}{Costos Constantes (\(S_{LP}\) Horizontal)}{
        Al expandirse la industria, el precio de los insumos \textbf{no varía} (industria pequeña frente al mercado de factores).\\[4pt]
        \(S_{LP}\) es perfectamente elástica al nivel \(P = \min CMe_{LP}\).
      }
    \end{column}
    \begin{column}{0.32\textwidth}
      \tarjeta{naranjaAcento}{Costos Crecientes (\(S_{LP}\) Positiva)}{
        Al entrar firmas, la mayor demanda de insumos específicos (ej. tierra o ingenieros especialistas) sube \((w, r)\) y eleva \(\min CMe_{LP}\).\\[4pt]
        \(S_{LP}\) tiene \textbf{pendiente positiva}.
      }
    \end{column}
    \begin{column}{0.32\textwidth}
      \tarjeta{verdeAcento}{Costos Decrecientes (\(S_{LP}\) Negativa)}{
        La expansión genera \textbf{externalidades pecuniarias positivas} (mejores redes de proveedores, infraestructura común) y baja \(\min CMe_{LP}\).\\[4pt]
        \(S_{LP}\) tiene \textbf{pendiente negativa}.
      }
    \end{column}
  \end{columns}

  \vspace{0.15cm}
  \recuadroclave{
    \small\textbf{Renta Económica vs. Cuasi-Renta:} A LP, si hay firmas con insumos escasos de menor costo, ellas obtienen \textbf{renta económica} (diferencia frente a la firma marginal), mientras la firma marginal tiene \(\pi = 0\). A CP, el \(EP = IT - CV = \pi + CF\) se llama \textbf{cuasi-renta}.
  }
  \mensajefinal{Un aumento en la demanda siempre sube \(Q^*\) en cualquier industria, pero solo sube \(P^*\) a LP en industrias de costos crecientes.}
\end{frame}

% ==============================================================================
% SLIDE 6: EXCEDENTES Y EFICIENCIA DE PARETO
% ==============================================================================
\begin{frame}{4. Excedentes del Consumidor, Productor y Bienestar Social}
  \begin{columns}[T]
    \begin{column}{0.53\textwidth}
      {\small\textbf{Maximización del Excedente Total (\(W = EC + EP\))}}
      \vspace{0.05cm}
      \begin{tikzpicture}[scale=0.78]
        \begin{axis}[
          axis lines=left,
          xlabel={Cantidad ($Q$)}, ylabel={Precio ($P$)},
          xmin=0, xmax=10, ymin=0, ymax=10,
          xtick={4.5}, xticklabels={$Q^*$},
          ytick={1.5, 5.1, 9.0}, yticklabels={$a$, $P^*$, $b$},
          width=8.2cm, height=5.2cm
        ]
          % Areas EC y EP
          \addplot[fill=azulClaroFondo, draw=none] coordinates {(0,5.1) (4.5,5.1) (0,9.0)} \closedcycle;
          \addplot[fill=verdeAcento!18, draw=none] coordinates {(0,1.5) (4.5,5.1) (0,5.1)} \closedcycle;
          % Demanda y Oferta
          \addplot[azulPrincipal, very thick, domain=0:9] {9 - 0.8667*x};
          \addplot[verdeAcento, very thick, domain=0:8] {1.5 + 0.8*x};
          % Lineas de equilibrio
          \draw[dashed, grisTexto] (axis cs:0, 5.1) -- (axis cs:4.5, 5.1) -- (axis cs:4.5, 0);
          \addplot[only marks, mark=*, mark size=2.6pt, rojoAcento] coordinates {(4.5, 5.1)};
          % Etiquetas
          \node[azulPrincipal, font=\small] at (axis cs:1.4, 6.4) {\textbf{EC}};
          \node[verdeAcento, font=\small] at (axis cs:1.4, 3.8) {\textbf{EP}};
          \node[azulPrincipal, font=\scriptsize, anchor=west] at (axis cs:7.0, 2.8) {$D$};
          \node[verdeAcento, font=\scriptsize, anchor=west] at (axis cs:7.0, 7.2) {$S$};
        \end{axis}
      \end{tikzpicture}
    \end{column}

    \begin{column}{0.44\textwidth}
      {\scriptsize\color{grisTexto}Excedente del Consumidor (\(EC\))}\\[1pt]
      {\large\color{azulPrincipal}\textbf{\(EC = \int_0^{Q^*} P^d(Q)\,dQ - P^* Q^*\)}}\\[6pt]
      {\scriptsize\color{grisTexto}Excedente del Productor (\(EP\))}\\[1pt]
      {\normalsize\color{verdeAcento}\textbf{\(EP = P^* Q^* - \int_0^{Q^*} CMg(Q)\,dQ = \pi + CF\)}}\\[6pt]
      {\scriptsize
        • En \(Q^*\), la disposición marginal a pagar (\(P^d\)) iguala al costo marginal de producir (\(CMg\)).\\
        • Cualquier desviación \(Q \neq Q^*\) destruye bienestar y genera una \textbf{Pérdida Irrecuperable de Eficiencia (PIE)}.\\
        • Es la expresión práctica del \textbf{Primer Teorema del Bienestar}.
      }
    \end{column}
  \end{columns}

  \mensajefinal{El equilibrio competitivo maximiza la suma \(EC + EP\), pero no garantiza equidad distributiva.}
\end{frame}

% ==============================================================================
% SLIDE 7: IMPUESTOS, INCIDENCIA Y PÉRDIDA DE EFICIENCIA (PIE)
% ==============================================================================
\begin{frame}{5. Impuestos Específicos, Incidencia y Peso Muerto (\(PIE\))}
  \begin{columns}[T]
    \begin{column}{0.53\textwidth}
      {\small\textbf{Impuesto Unitario \(t\): Brecha \(P_c - P_p = t\)}}
      \vspace{0.05cm}
      \begin{tikzpicture}[scale=0.78]
        \begin{axis}[
          axis lines=left,
          xlabel={Cantidad ($Q$)}, ylabel={Precio ($P$)},
          xmin=0, xmax=10, ymin=0, ymax=10,
          xtick={3.0, 4.5}, xticklabels={$Q_t$, $Q^*$},
          ytick={3.9, 5.1, 6.4}, yticklabels={$P_p$, $P^*$, $P_c$},
          width=8.2cm, height=5.2cm
        ]
          % Areas: Recaudacion y PIE
          \addplot[fill=naranjaAcento!20, draw=none] coordinates {(0,3.9) (3.0,3.9) (3.0,6.4) (0,6.4)} \closedcycle;
          \addplot[fill=rojoAcento!30, draw=none] coordinates {(3.0,3.9) (4.5,5.1) (3.0,6.4)} \closedcycle;
          % Demanda y Oferta
          \addplot[azulPrincipal, very thick, domain=0:9] {9 - 0.8667*x};
          \addplot[verdeAcento, very thick, domain=0:8] {1.5 + 0.8*x};
          \addplot[verdeAcento, thick, dashed, domain=0:6.5] {4.0 + 0.8*x};
          % Lineas punteadas
          \draw[dotted, grisTexto] (axis cs:0, 6.4) -- (axis cs:3.0, 6.4) -- (axis cs:3.0, 0);
          \draw[dotted, grisTexto] (axis cs:0, 3.9) -- (axis cs:3.0, 3.9);
          \draw[dotted, grisTexto] (axis cs:0, 5.1) -- (axis cs:4.5, 5.1) -- (axis cs:4.5, 0);
          % Etiquetas
          \node[naranjaAcento, font=\scriptsize] at (axis cs:1.5, 5.15) {\textbf{Recaud.}};
          \node[rojoAcento, font=\scriptsize] at (axis cs:3.45, 5.15) {\textbf{PIE}};
        \end{axis}
      \end{tikzpicture}
    \end{column}

    \begin{column}{0.44\textwidth}
      \indicador{Incidencia Económica}{naranjaAcento}{\(\frac{\Delta P_c}{|\Delta P_p|} = \frac{\varepsilon_s}{|\varepsilon_d|}\)}
      {\scriptsize
        • La \textbf{incidencia legal} (quién paga en ventanilla) es irrelevante.\\
        • \textbf{Regla de oro:} El lado \textbf{más inelástico} del mercado absorbe la mayor parte del impuesto.\\[4pt]
      }
      {\scriptsize\color{grisTexto}Pérdida Irrecuperable (Triángulo de Harberger)}\\[1pt]
      {\normalsize\color{rojoAcento}\textbf{\(PIE = \frac{1}{2} t (Q^* - Q_t) \approx \frac{1}{2}\frac{t^2 Q^*}{P^*} \frac{\varepsilon_s |\varepsilon_d|}{\varepsilon_s + |\varepsilon_d|}\)}}\\[5pt]
      {\scriptsize
        • La \(PIE\) crece con el \textbf{cuadrado} de la tasa impositiva (\(t^2\)) y es mayor cuanto más \textbf{elásticas} son \(D\) y \(S\).\\
        • Un subsidio \(s\) también genera \(PIE = \frac{1}{2} s \Delta Q\).
      }
    \end{column}
  \end{columns}

  \mensajefinal{Duplicar la tasa de un impuesto cuadruplica la pérdida irrecuperable de eficiencia en el mercado.}
\end{frame}

% ==============================================================================
% SLIDE 8: CONTROLES DE PRECIOS, CUOTAS Y ARANCELES
% ==============================================================================
\begin{frame}{6. Controles de Precios, Cuotas y Apertura Comercial}
  \vspace{0.1cm}
  \begin{columns}[T]
    \begin{column}{0.32\textwidth}
      \tarjeta{rojoAcento}{Precio Máximo (Tope)}{
        Efectivo solo si \(P_{\max} < P^*\).\\[4pt]
        • Genera \textbf{escasez} (\(Q^d > Q^s\)) y racionamiento/colas.\\
        • \textbf{Siempre reduce el \(EP\)}; puede subir o bajar el \(EC\), y siempre reduce el bienestar total (\(PIE > 0\)).
      }
    \end{column}
    \begin{column}{0.32\textwidth}
      \tarjeta{naranjaAcento}{Precio Mínimo (Piso)}{
        Efectivo solo si \(P_{\min} > P^*\).\\[4pt]
        • Genera \textbf{excedente de oferta} (\(Q^s > Q^d\)).\\
        • Ej. salario mínimo en mercado competitivo: causa desempleo si \(w_{\min} > w^*\).
      }
    \end{column}
    \begin{column}{0.32\textwidth}
      \tarjeta{azulPrincipal}{Economía Abierta y Arancel}{
        Si \(P_w < P_{aut}\): el país \textbf{importa} y gana bienestar neto.\\[4pt]
        • Un \textbf{arancel \(t\)} eleva \(P = P_w + t\), reduce importaciones y crea dos triángulos de \(PIE\) (distorsión de producción y consumo).
      }
    \end{column}
  \end{columns}

  \vspace{0.15cm}
  \recuadroclave{
    \small\textbf{Regla en Controles de Precios:} Cuando un precio tope causa escasez, el volumen transado lo determina el lado corto (\(Q = Q^s(P_{\max})\)) y la \(PIE\) depende de cómo se raciona el bien.
  }
  \mensajefinal{El libre comercio siempre eleva el excedente agregado frente a la autarquía, aunque genera ganadores y perdedores internos.}
\end{frame}

% ==============================================================================
% SLIDE 9: AD-VALOREM VS ESPECÍFICO, CUOTAS Y MODELO DE LA TELARAÑA
% ==============================================================================
\begin{frame}{7. Ad-Valorem, Cuotas de Producción y Modelo de la Telaraña}
  \vspace{0.05cm}
  \begin{columns}[T]
    \begin{column}{0.32\textwidth}
      \tarjeta{azulPrincipal}{Ad-Valorem vs. Específico}{
        • \textbf{Específico (\(t\)):} \(P_c = P_p + t\) (traslación paralela).\\
        • \textbf{Ad-Valorem (\(\tau\)):} \(P_c = P_p(1+\tau)\) (rotación de la curva).\\
        • En \textbf{competencia perfecta}, ambos generan la \textbf{misma \(PIE\)} para igual recaudación.
      }
    \end{column}
    \begin{column}{0.32\textwidth}
      \tarjeta{naranjaAcento}{Arancel vs. Cuota de Importación}{
        Para igual volumen restringido, el arancel y la cuota generan igual \(P\), \(Q\) y \(PIE\).\\[2pt]
        \textbf{Diferencia clave:} En el arancel, el rectángulo \((P_d - P_w)\cdot M\) es \textbf{recaudación fiscal}; en la cuota, se convierte en \textbf{renta de cuota} para quien recibe la licencia.
      }
    \end{column}
    \begin{column}{0.32\textwidth}
      \tarjeta{verdeAcento}{Modelo de la Telaraña (Cobweb)}{
        Oferta con rezago: \(Q_t^s = S(P_{t-1})\), \(Q_t^d = D(P_t)\):\\
        • \textbf{Convergente:} \(|dP/dQ_S| > |dP/dQ_D|\) (oferta más inelástica que demanda).\\
        • \textbf{Divergente:} \(|dP/dQ_S| < |dP/dQ_D|\).\\
        • \textbf{Ciclo perpetuo:} pendientes iguales.
      }
    \end{column}
  \end{columns}

  \vspace{0.1cm}
  \recuadroclave{
    \small\textbf{Subsidio en Equilibrio Competitivo:} \(P_p - P_c = s\). Eleva \(Q\) por encima de \(Q^*\) (sobreproducción) y genera \(PIE = \frac{1}{2} s (Q_s - Q^*)\) porque el gasto estatal supera la ganancia de \(EC + EP\).
  }
  \mensajefinal{En competencia perfecta los impuestos ad-valorem y específicos son equivalentes en bienestar; en monopolio el ad-valorem domina.}
\end{frame}

% ==============================================================================
% SLIDE 10: CONCLUSIONES Y FORMULARIO RÁPIDO MÓDULO III
% ==============================================================================
\begin{frame}{Conclusiones y Formulario Rápido — Módulo III}
  \vspace{0.05cm}
  \recuadroclave{
    \small\textbf{Regla de Oro:} La competencia perfecta logra eficiencia paretiana en el punto donde \(P = CMg\) (asignativa) y a largo plazo en \(P = \min CMe_{LP}\) (productiva).
  }

  \vspace{0.05cm}
  \begin{columns}[T]
    \begin{column}{0.32\textwidth}
      \tarjeta{verdeAcento}{Fórmulas Imprescindibles}{
        • Equilibrio LP: \(P = CMg_{LP} = \min CMe_{LP}\)\\
        • Incidencia: \(\frac{\Delta P_c}{|\Delta P_p|} = \frac{\varepsilon_s}{|\varepsilon_d|}\)\\
        • \(PIE = \frac{1}{2} t \Delta Q\)
      }
    \end{column}
    \begin{column}{0.32\textwidth}
      \tarjeta{azulPrincipal}{Trampas de Examen}{
        • Creer que \(\pi=0\) a LP implica cero ganancia contable (incluye beneficio normal).\\
        • Confundir \(S_{LP} > 0\) con falla de mercado (es por costos crecientes).\\
        • Olvidar que \(CF\) no mueve \(S_{CP}\).
      }
    \end{column}
    \begin{column}{0.32\textwidth}
      \tarjeta{rojoAcento}{Siguientes Pasos}{
        • Practicar Bloques B y C del Módulo 03 en el Simulador.\\
        • Enlazar con el Módulo IV (Equilibrio General y Bienestar).
      }
    \end{column}
  \end{columns}

  \vspace{0.12cm}
  \begin{center}
    {\color{azulPrincipal}\Large\textbf{¡Éxito!}} \quad {\color{grisTexto}Módulo III · Competencia Perfecta · PEU SUNASS 2026}
  \end{center}
\end{frame}

\end{document}
